% Generated by publication/build_sections.py; do not edit by hand.
\paragraph{Informal verdict.}
\textbf{A --- complete and apparently correct.}

\paragraph{Lean coverage.}
The complete Q655Stage7 closure proves the maximum n\textasciicircum{}2-1 over every field and n>=2, the exact n=2 iff noncommutation classification, and for every n>=3 the symmetric iff classification in which one member is scalar plus nonzero rank one and the other satisfies exactly the split, no-root quadratic (including inseparable characteristic 2), or repeated-root square-zero conditions. Stage5 proves the scalar-rank inequality without an extra premise.

\subsection{Audited informal answer}
Complete solution
\begin{enumerate}
\item Verified statement
\end{enumerate}
The canonical PDF states, in the section Algèbre linéaire et bilinéaire:

Q655 (following R626). Let \(K\) be an arbitrary field and let \(E\) be a finite-dimensional \(K\)-vector space of dimension \(n\ge 2\). For \(f\in\mathcal L(E)\), denote by \(C(f)\) its commutant. When neither \(f\) nor \(g\) is a scalar endomorphism, what is the maximum possible dimension of \(C(f)+C(g)\), and which pairs \((f,g)\) attain equality?

There are no subparts, footnotes, or ambiguities attached to Q655. Luc Pommeret
The supplied previous attempt correctly isolated the sole missing implication: whether arbitrary-dimensional extremality forces one member to differ from a scalar by a rank-one map. Fichier markdown.md collé The centralizer-sum inequality proved below establishes exactly that implication.

\begin{enumerate}
\item Final answer
\end{enumerate}
Let \(A=\operatorname{End}_K(E)\).
Theorem 2.1 --- Maximum and exhaustive equality classification
For every field \(K\) and every \(n\ge 2\),

\[
\boxed{
\max_{\substack{f,g\in A\\ f,g\ \mathrm{non\text{-}scalar}}}
\dim_K\bigl(C(f)+C(g)\bigr)=n^2-1.
}
\]

The equality pairs are as follows.
Dimension \(n=2\)
Every non-scalar endomorphism has a two-dimensional commutant, and

\[
\boxed{
\dim(C(f)+C(g))=3
\iff C(f)\ne C(g)
\iff fg\ne gf.
}
\]

This includes pairs whose minimal polynomials are irreducible quadratics over \(K\).
Dimensions \(n\ge 3\)
Equality holds if and only if, after possibly interchanging \(f\) and \(g\),

\[
g=\lambda I+u\otimes\varphi,
\qquad
u\ne0,\quad \varphi\ne0,
\tag{2.1}
\]

and \(f\) has quadratic minimal polynomial and satisfies one of the following conditions.
I. Split quadratic case
Suppose

\[
\mu_f(X)=(X-a)(X-b),\qquad a\ne b,
\]

and write

\[
E=E_a\oplus E_b,\qquad
u=u_a+u_b,\qquad
\varphi=\varphi_a+\varphi_b,
\]

where \(\varphi_i=\varphi|_{E_i}\). For each \(i\in\{a,b\}\),

\[
\begin{cases}
u_i\ne0\ \text{and}\ \varphi_i\ne0,
   &\dim E_i\ge2,\\[2mm]
(u_i,\varphi_i)\ne(0,0),
   &\dim E_i=1.
\end{cases}
\tag{2.2}
\]

II. Irreducible quadratic case
Suppose \(\mu_f\) is irreducible of degree \(2\) over \(K\). Then every choice

\[
u\ne0,\qquad \varphi\ne0
\tag{2.3}
\]

gives equality. Necessarily \(n\) is even.
This includes inseparable irreducible quadratics in characteristic \(2\); no separability assumption is needed.
III. Repeated-root quadratic case
Suppose

\[
f=aI+N,\qquad N^2=0,\qquad N\ne0,
\]

and put \(k=\operatorname{rank}N\). Equality holds exactly in one of the following mutually exclusive situations:

\[
Nu\ne0,\qquad \varphi N\ne0;
\tag{2.4}
\]

or

\[
Nu=0,\qquad
\varphi N\ne0,\qquad
k=1,\qquad
u\notin\operatorname{im}N;
\tag{2.5}
\]

or

\[
Nu\ne0,\qquad
\varphi N=0,\qquad
k=1,\qquad
\varphi|_{\ker N}\ne0.
\tag{2.6}
\]

If \(Nu=0\) and \(\varphi N=0\), equality is impossible.
In every equality case with \(n\ge3\),

\[
\dim C(g)=n^2-2n+2
\tag{2.7}
\]

and

\[
\boxed{
\dim\bigl(C(f)\cap C(g)\bigr)
=
\dim C(f)-2n+3.
}
\tag{2.8}
\]

Thus every commutant and intersection dimension is explicitly determined.

\begin{enumerate}
\item The decisive centralizer-sum inequality
\end{enumerate}
For an endomorphism \(a\), define its geometric scalar rank

\[
\rho(a)
:=
\min_{\lambda\in\overline K}
\operatorname{rank}_{\overline K}(a-\lambda I),
\tag{3.1}
\]

where \(\overline K\) is an algebraic closure of \(K\).
Equivalently,

\[
\rho(a)
=
n-\max_{\lambda\in\overline K}
\dim_{\overline K}\ker(a-\lambda I).
\tag{3.2}
\]

The missing arbitrary-dimensional implication follows from the following stronger result.
Theorem 3.1 --- Scalar-rank bound
For arbitrary \(f,g\in\operatorname{End}_K(E)\),

\[
\boxed{
\dim_K\bigl(C(f)+C(g)\bigr)
\le
n^2-\min\{\rho(f),\rho(g)\}.
}
\tag{3.3}
\]

In particular, if neither endomorphism differs from a scalar by rank at most \(r-1\), then

\[
\dim(C(f)+C(g))\le n^2-r.
\tag{3.4}
\]

Thus, for \(n\ge3\), if neither member is a scalar plus a rank-one map,

\[
\boxed{\dim(C(f)+C(g))\le n^2-2.}
\tag{3.5}
\]

This immediately eliminates every proposed counterexample in dimension \(6\), in characteristic \(2\), and in all larger dimensions.

3.1. Scalar extension
For every \(a\in\operatorname{End}_K(E)\),

\[
C_{\overline K}(a_{\overline K})
=
C_K(a)\otimes_K\overline K,
\]

because both sides are the kernel after scalar extension of the same commutator map. Consequently, all relevant centralizer, intersection, and sum dimensions are unchanged by scalar extension.
We may therefore prove Theorem 3.1 over an algebraically closed field.
Put

\[
r=\rho(f),\qquad s=\rho(g),
\]

and interchange \(f,g\) if necessary so that

\[
r\le s.
\]

Choose scalars \(\lambda,\mu\) such that

\[
F=f-\lambda I,\qquad G=g-\mu I,
\qquad
\operatorname{rank}F=r,\quad
\operatorname{rank}G=s.
\tag{3.6}
\]

Scalar shifts do not change commutants.

3.2. A centralizer estimate when \(r+s>n\)
We first record a standard estimate.
Lemma 3.2
If \(a\) is an endomorphism of an \(n\)-dimensional vector space over an algebraically closed field, and

\[
e(a)=\max_\alpha\dim\ker(a-\alpha I),
\]

then

\[
\dim C(a)\le n\,e(a).
\tag{3.7}
\]

Proof
Let \(c_{\alpha,j}\) denote the column lengths of the Jordan partition belonging to the eigenvalue \(\alpha\). Counting homomorphisms between Jordan blocks gives

\[
\dim C(a)=\sum_{\alpha,j}c_{\alpha,j}^2.
\]

Every column length is at most \(e(a)\), while the sum of all column lengths is \(n\). Hence

\[
\dim C(a)
\le
e(a)\sum_{\alpha,j}c_{\alpha,j}
=
ne(a).
\qquad\square
\]

Here \(e(f)=n-r\) and \(e(g)=n-s\). If \(r+s>n\), then, since both centralizers contain \(KI\),

\[
\begin{aligned}
\dim(C(f)+C(g))
&\le \dim C(f)+\dim C(g)-1\\
&\le n(n-r)+n(n-s)-1\\
&=n(2n-r-s)-1\\
&\le n(n-1)-1\\
&\le n^2-r.
\end{aligned}
\tag{3.8}
\]

Thus Theorem 3.1 is proved in this case.

3.3. Rank factorizations and their centralizers
Assume henceforth that

\[
r+s\le n.
\tag{3.9}
\]

We need two exact factorization lemmas.
Lemma 3.3 --- Centralizer of a rank factorization
Let \(U:L\to E\) be injective, let \(V:E\to L\) be surjective, put

\[
F=UV,\qquad A_0=VU,
\]

and let \(\dim L=r\). Then there is an exact sequence

\[
0\longrightarrow
\operatorname{Hom}(E/\operatorname{im}U,\ker V)
\longrightarrow
C(F)
\longrightarrow
C(A_0)
\longrightarrow0.
\tag{3.10}
\]

Consequently,

\[
\boxed{
\dim C(F)=(n-r)^2+\dim C(A_0).
}
\tag{3.11}
\]

Proof
If \(X\in C(F)\), then \(X\) preserves \(\operatorname{im}F=\operatorname{im}U\). Hence there is a unique \(P\in\operatorname{End}(L)\) such that

\[
XU=UP.
\]

The equality \(XF=FX\), together with the injectivity of \(U\), then gives

\[
PV=VX.
\]

Applying \(V\) to \(XU=UP\) shows that

\[
A_0P=PA_0.
\]

Thus \(X\mapsto P\) defines a map \(C(F)\to C(A_0)\).
Conversely, let \(P\in C(A_0)\). Choose a right inverse \(S:L\to E\) of \(V\), and set

\[
X_0=SPV.
\]

Then \(VX_0=PV\). Moreover,

\[
Y:=UP-X_0U
\]

has image in \(\ker V\), since

\[
VY=A_0P-PA_0=0.
\]

Define \(Z\) on \(\operatorname{im}U\) by

\[
Z(U\ell)=Y\ell,
\]

and extend \(Z\) to \(E\) with image contained in \(\ker V\). Then

\[
X=X_0+Z
\]

satisfies \(XU=UP\) and \(VX=PV\), hence commutes with \(F=UV\).
The kernel consists precisely of the maps vanishing on \(\operatorname{im}U\) and having image in \(\ker V\). Its dimension is

\[
(n-r)(n-r).
\qquad\square
\]

Apply this to rank factorizations

\[
F=U_1V_1,\qquad G=U_2V_2,
\tag{3.12}
\]

where \(U_1,U_2\) are injective and \(V_1,V_2\) are surjective. Put

\[
A_0=V_1U_1\in M_r(K),\qquad
D_0=V_2U_2\in M_s(K).
\tag{3.13}
\]

Then

\[
\dim C(f)=(n-r)^2+\dim C(A_0),
\tag{3.14}
\]


\[
\dim C(g)=(n-s)^2+\dim C(D_0).
\tag{3.15}
\]


3.4. Generic relative position
Conjugate \(G\) relative to \(F\). The function

\[
h\longmapsto
\dim\bigl(C(F)\cap C(hGh^{-1})\bigr)
\]

is the nullity of a matrix whose coefficients are regular functions of \(h\in\operatorname{GL}(E)\). Its minimum is therefore attained on a nonempty Zariski-open subset.
There are also nonempty open subsets on which

\[
\operatorname{im}F\cap h(\operatorname{im}G)=0
\]

and

\[
\dim\bigl(\ker F\cap h(\ker G)\bigr)=n-r-s.
\]

Because \(\operatorname{GL}(E)\) is irreducible, these open sets intersect. Hence the minimum possible common-centralizer dimension occurs for a relative position in which

\[
U:K^r\oplus K^s\longrightarrow E,
\qquad
U(x,y)=U_1x+U_2y
\]

is injective and

\[
V:E\longrightarrow K^r\oplus K^s,
\qquad
V(v)=(V_1v,V_2v)
\]

is surjective.
Write

\[
VU=
\begin{pmatrix}
A_0&B\\
C&D_0
\end{pmatrix}.
\tag{3.16}
\]

Define

\[
\mathcal E(B,C)=
\left\{
(P,Q)\in C(A_0)\oplus C(D_0):
\begin{array}{l}
PB=BQ,\\
QC=CP
\end{array}
\right\}.
\tag{3.17}
\]

Exactly as in Lemma 3.3, one obtains an exact sequence

\[
0\longrightarrow
\operatorname{Hom}(E/\operatorname{im}U,\ker V)
\longrightarrow
C(F)\cap C(G)
\longrightarrow
\mathcal E(B,C)
\longrightarrow0.
\tag{3.18}
\]

Since both \(E/\operatorname{im}U\) and \(\ker V\) have dimension \(n-r-s\),

\[
\dim\bigl(C(F)\cap C(G)\bigr)
=
(n-r-s)^2+\dim\mathcal E(B,C)
\tag{3.19}
\]

in this minimum configuration.

3.5. The two-arrow estimate
Lemma 3.4
For \(r\le s\),

\[
\boxed{
\dim\mathcal E(B,C)
\ge
\dim C(A_0)+\dim C(D_0)-2rs+r.
}
\tag{3.20}
\]

Proof
For variable pairs

\[
(B,C)\in
M_{r,s}(K)\oplus M_{s,r}(K),
\]

consider the linear map

\[
\Phi_{B,C}:C(A_0)\oplus C(D_0)
\longrightarrow
M_{r,s}(K)\oplus M_{s,r}(K),
\]


\[
\Phi_{B,C}(P,Q)
=
(PB-BQ,\ QC-CP).
\tag{3.21}
\]

Its kernel is \(\mathcal E(B,C)\).
Let

\[
\Psi(B,C)
\]

be the \(r\) non-leading coefficients of the characteristic polynomial of \(BC\in M_r(K)\). There is a nonempty open subset on which \(d\Psi\) has rank \(r\). Indeed, take

\[
B=\begin{pmatrix}I_r&0\end{pmatrix},
\qquad
C=\begin{pmatrix}T\\0\end{pmatrix},
\]

where \(T\) is a companion matrix. Varying the companion coefficients gives \(r\) independent variations of the characteristic polynomial.
The maximum-rank locus of \(\Phi_{B,C}\) is also nonempty and open. Choose \((B,C)\) in the intersection of these two open sets.
For \(P\in C(A_0)\) and \(Q\in C(D_0)\), the infinitesimal variation

\[
(\delta B,\delta C)
=
(PB-BQ,\ QC-CP)
\]

gives

\[
\delta(BC)
=
(PB-BQ)C+B(QC-CP)
=
PBC-BCP.
\]

Thus \(BC\) varies by an infinitesimal conjugation. Its characteristic polynomial does not change, so

\[
d\Psi\bigl(\Phi_{B,C}(P,Q)\bigr)=0.
\]

Therefore

\[
\operatorname{rank}\Phi_{B,C}
\le 2rs-r.
\]

Because this point was chosen in the maximum-rank locus, the same bound holds for every \(B,C\). Hence

\[
\begin{aligned}
\dim\mathcal E(B,C)
&=\dim C(A_0)+\dim C(D_0)
  -\operatorname{rank}\Phi_{B,C}\\
&\ge
\dim C(A_0)+\dim C(D_0)-2rs+r.
\end{aligned}
\]

No division is used, so the argument applies unchanged in characteristic \(2\). \(\square\)

3.6. Completion of the scalar-rank bound
Combining (3.19) and Lemma 3.4 gives, for the minimum common-centralizer dimension,

\[
\begin{aligned}
\dim(C(F)\cap C(G))
\ge{}&(n-r-s)^2\\
&+\dim C(A_0)+\dim C(D_0)-2rs+r.
\end{aligned}
\tag{3.22}
\]

Every other relative position has common-centralizer dimension at least this large. Hence, using (3.14) and (3.15),

\[
\begin{aligned}
\dim(C(f)+C(g))
&=\dim C(f)+\dim C(g)-\dim(C(f)\cap C(g))\\
&\le
(n-r)^2+(n-s)^2\\
&\quad -(n-r-s)^2+2rs-r\\
&=n^2-r.
\end{aligned}
\]

Since \(r=\min\{\rho(f),\rho(g)\}\), this proves Theorem 3.1. \(\square\)

\begin{enumerate}
\item Extremality forces a rank-one member
\end{enumerate}
Suppose now that \(f,g\) are non-scalar and

\[
\dim(C(f)+C(g))=n^2-1.
\tag{4.1}
\]

Theorem 3.1 gives

\[
n^2-1
\le
n^2-\min\{\rho(f),\rho(g)\},
\]

so

\[
\min\{\rho(f),\rho(g)\}=1.
\tag{4.2}
\]

Thus, over \(\overline K\), one of the two endomorphisms differs from a scalar by rank one.
For \(n\ge3\), this property descends to \(K\).
Lemma 4.1 --- Descent of the scalar
Let \(a\in M_n(K)\), \(n\ge3\), and suppose

\[
\operatorname{rank}_{\overline K}(a-\lambda I)=1
\]

for some \(\lambda\in\overline K\). Then \(\lambda\in K\).
Proof
If \(a\) is diagonal, rank one implies that at least \(n-1\) diagonal entries equal \(\lambda\), hence \(\lambda\in K\).
Otherwise choose an off-diagonal entry \(a_{ij}\ne0\), and choose \(k\notin\{i,j\}\). The \(2\times2\) minor of \(a-\lambda I\) with rows \(i,k\) and columns \(j,k\) vanishes:

\[
a_{ij}(a_{kk}-\lambda)-a_{ik}a_{kj}=0.
\]

Therefore

\[
\lambda
=
a_{kk}-\frac{a_{ik}a_{kj}}{a_{ij}}
\in K.
\qquad\square
\]

Consequently, for \(n\ge3\), after interchanging the pair,

\[
g=\lambda I+r,
\qquad
\operatorname{rank}r=1.
\tag{4.3}
\]

This proves the formerly missing necessity statement in every dimension and every characteristic.

\begin{enumerate}
\item Exact reduction when one member has rank one
\end{enumerate}
Write

\[
r=u\otimes\varphi,
\qquad u\ne0,\quad\varphi\ne0.
\]

Since scalar shifts do not affect commutants, \(C(g)=C(r)\).
Lemma 5.1 --- Rank-one centralizer

\[
\boxed{
\dim C(r)=n^2-2n+2.
}
\tag{5.1}
\]

Proof
Consider

\[
S:\operatorname{End}(E)\longrightarrow E\oplus E^*,
\qquad
S(X)=(Xu,\varphi X).
\]

Its kernel consists of maps \(E/Ku\to\ker\varphi\), hence has dimension \((n-1)^2\).
The equality \(Xr=rX\) is equivalent to the existence of a scalar \(c\) such that

\[
Xu=cu,\qquad \varphi X=c\varphi.
\]

Thus

\[
C(r)=S^{-1}\bigl(K(u,\varphi)\bigr).
\]

Since \(S(I)=(u,\varphi)\),

\[
\dim C(r)=(n-1)^2+1=n^2-2n+2.
\qquad\square
\]

Define

\[
T_{f,u,\varphi}:C(f)\longrightarrow E\oplus E^*,
\qquad
T(X)=(Xu,\varphi X).
\tag{5.2}
\]

Lemma 5.2 --- Exact sum formula

\[
\boxed{
\dim(C(f)+C(r))
=
n^2-2n+1+\operatorname{rank}T_{f,u,\varphi}.
}
\tag{5.3}
\]

Proof
For \(X\in C(f)\),

\[
[r,X]
=
u\otimes(\varphi X)-(Xu)\otimes\varphi.
\]

The linear map

\[
L:E\oplus E^*\to\operatorname{End}(E),
\qquad
L(v,\alpha)=u\otimes\alpha-v\otimes\varphi
\]

has kernel \(K(u,\varphi)\). Since \(T(I)=(u,\varphi)\), this kernel line lies in \(\operatorname{im}T\). Hence

\[
\operatorname{rank}\bigl(\operatorname{ad}_r|_{C(f)}\bigr)
=
\operatorname{rank}T-1.
\]

Its kernel is \(C(f)\cap C(r)\). Consequently,

\[
\begin{aligned}
\dim(C(f)+C(r))
&=\dim C(r)
 +\operatorname{rank}\bigl(\operatorname{ad}_r|_{C(f)}\bigr)\\
&=n^2-2n+2+\operatorname{rank}T-1.
\end{aligned}
\]

\(\square\)
Therefore,

\[
\dim(C(f)+C(r))=n^2-1
\iff
\operatorname{rank}T=2n-2.
\tag{5.4}
\]


\begin{enumerate}
\item Why \(f\) must be quadratic
\end{enumerate}
Lemma 6.1
Let \(d=\deg\mu_f\). Then

\[
\boxed{
\operatorname{rank}T_{f,u,\varphi}\le 2n-d.
}
\tag{6.1}
\]

Proof
Scalar extension preserves both sides, so assume that \(K\) is algebraically closed.
For \(0\le i<d\), define on \(E\oplus E^*\)

\[
q_i(v,\alpha)=\alpha(f^iv).
\]

For \(X\in C(f)\), consider the tangent vector

\[
(Xv,-\alpha X).
\]

Then

\[
\begin{aligned}
dq_i{}_{(v,\alpha)}(Xv,-\alpha X)
&=-(\alpha X)(f^iv)+\alpha(f^iXv)\\
&=0.
\end{aligned}
\tag{6.2}
\]

There is \(v_0\in E\) such that

\[
v_0,fv_0,\ldots,f^{d-1}v_0
\]

are linearly independent. At \((v_0,0)\), the \(d\) differentials \(dq_i\) are independent in the \(E^*\)-directions.
The locus where these differentials are independent is nonempty and open. The locus where the map

\[
X\longmapsto(Xv,-\alpha X)
\]

has maximum rank is also nonempty and open. Their intersection is nonempty. At a point in the intersection, the image is annihilated by \(d\) independent covectors, and therefore has dimension at most \(2n-d\). This bounds its maximum rank, hence its rank everywhere.
Changing the sign in the second component does not change rank, so the same estimate holds for \(T\). \(\square\)
Under equality,

\[
2n-2=\operatorname{rank}T\le2n-\deg\mu_f.
\]

Since \(f\) is non-scalar,

\[
\boxed{\deg\mu_f=2.}
\tag{6.3}
\]

It remains only to analyze the three kinds of quadratic polynomial.

\begin{enumerate}
\item Classification of the quadratic cases
\end{enumerate}
7.1. Split quadratic
Suppose

\[
\mu_f=(X-a)(X-b),\qquad a\ne b.
\]

Then

\[
E=E_a\oplus E_b,\qquad
C(f)=\operatorname{End}(E_a)\oplus\operatorname{End}(E_b).
\]

For an \(m\)-dimensional space \(W\), define

\[
F_{x,\xi}:\operatorname{End}(W)\to W\oplus W^*,
\qquad
A\mapsto(Ax,\xi A).
\]

A direct kernel calculation gives

\[
\operatorname{rank}F_{x,\xi}
=
\begin{cases}
0,&x=0,\ \xi=0,\\
m,&\text{exactly one of }x,\xi\text{ is nonzero},\\
2m-1,&x\ne0,\ \xi\ne0.
\end{cases}
\tag{7.1}
\]

Indeed, in the last case the kernel is

\[
\operatorname{Hom}(W/Kx,\ker\xi),
\]

of dimension \((m-1)^2\).
The map \(T\) is the direct sum of the maps associated with

\[
(u_a,\varphi_a)
\quad\text{and}\quad
(u_b,\varphi_b).
\]

Its required rank is

\[
2\dim E_a-1+2\dim E_b-1=2n-2.
\]

This gives exactly the conditions in (2.2).
In this case, if

\[
p=\dim E_a,\qquad q=\dim E_b,
\]

then

\[
\dim C(f)=p^2+q^2
\]

and, under equality,

\[
\dim(C(f)\cap C(g))=p^2+q^2-2n+3.
\tag{7.2}
\]


7.2. Irreducible quadratic
Suppose \(\mu_f=q\) is irreducible of degree \(2\), and let

\[
L=K[X]/(q).
\]

Then \(E\) is an \(L\)-vector space of dimension

\[
m=\frac n2,
\]

and

\[
C(f)=\operatorname{End}_L(E).
\]

Choose a nonzero \(K\)-linear functional

\[
\ell:L\to K.
\]

Composition with \(\ell\) gives a \(K\)-linear isomorphism

\[
\operatorname{Hom}_L(E,L)\longrightarrow E^*,
\qquad
\psi\mapsto\ell\circ\psi.
\tag{7.3}
\]

Thus write

\[
\varphi=\ell\circ\psi
\]

with \(\psi\ne0\).
The kernel of

\[
\operatorname{End}_L(E)\longrightarrow
E\oplus\operatorname{Hom}_L(E,L),
\qquad
A\mapsto(Au,\psi A)
\]

is

\[
\operatorname{Hom}_L(E/Lu,\ker\psi),
\]

which has \(L\)-dimension \((m-1)^2\). Hence the image has \(L\)-dimension

\[
m^2-(m-1)^2=2m-1.
\]

Its \(K\)-dimension is therefore

\[
2(2m-1)=4m-2=2n-2.
\]

So every nonzero \(u,\varphi\) gives equality.
Furthermore,

\[
\dim_K C(f)=2m^2=\frac{n^2}{2},
\]

and

\[
\dim(C(f)\cap C(g))
=
\frac{n^2}{2}-2n+3.
\tag{7.4}
\]

Nothing in this argument assumes that \(L/K\) is separable.

7.3. Repeated root
Suppose

\[
f=aI+N,\qquad N^2=0,\quad N\ne0.
\]

Put

\[
k=\operatorname{rank}N,\qquad m=n-2k.
\]

Choose a decomposition

\[
E=U\oplus W\oplus V,
\]

where

\[
U=\operatorname{im}N,\qquad
\ker N=U\oplus W,
\]

and \(N:V\to U\) is the identity under a fixed identification. Thus

\[
N=
\begin{pmatrix}
0&0&I\\
0&0&0\\
0&0&0
\end{pmatrix}.
\]

The centralizer consists exactly of matrices

\[
X=
\begin{pmatrix}
A&B&C\\
0&E_0&F\\
0&0&A
\end{pmatrix}.
\tag{7.5}
\]

Write

\[
u=(x,y,z),\qquad
\varphi=(\alpha,\beta,\gamma).
\]

Then

\[
Nu\ne0\iff z\ne0,
\qquad
\varphi N\ne0\iff\alpha\ne0.
\tag{7.6}
\]

When \(z=0\),

\[
u\notin\operatorname{im}N\iff y\ne0,
\]

and, when \(\alpha=0\),

\[
\varphi|_{\ker N}\ne0\iff\beta\ne0.
\]

The annihilator of \(\operatorname{im}T\) in

\[
(E\oplus E^*)^*\simeq E^*\oplus E
\]

consists of tuples

\[
(p,q,\rho;\,a_0,b,c)
\]

satisfying

\[
\begin{aligned}
x\otimes p+z\otimes\rho+a_0\otimes\alpha+c\otimes\gamma&=0,\\
y\otimes p+b\otimes\alpha&=0,\\
z\otimes p+c\otimes\alpha&=0,\\
y\otimes q+b\otimes\beta&=0,\\
z\otimes q+c\otimes\beta&=0.
\end{aligned}
\tag{7.7}
\]

Solving these rank-one tensor equations gives

\[
\begin{array}{c|c|c}
z&\alpha&
\dim(\operatorname{im}T)^\perp\\ \hline
z\ne0&\alpha\ne0&2\\
z=0&\alpha\ne0,\ y\ne0&k+1\\
z=0&\alpha\ne0,\ y=0&k+m+1\\
z\ne0&\alpha=0,\ \beta\ne0&k+1\\
z\ne0&\alpha=0,\ \beta=0&k+m+1\\
z=0&\alpha=0&>2.
\end{array}
\tag{7.8}
\]

For example, if \(z,\alpha\ne0\), the third equation gives

\[
p=t\alpha,\qquad c=-tz.
\]

The second and fifth give

\[
b=-ty,\qquad q=t\beta.
\]

The first equation introduces one further scalar parameter. Hence the annihilator is two-dimensional.
If \(z=0,\alpha\ne0,y\ne0\), then \(\rho\in V^*\) is arbitrary and the remaining variables depend on one scalar, giving dimension \(k+1\). The other rows follow identically or by duality.
Since

\[
\operatorname{rank}T
=
2n-\dim(\operatorname{im}T)^\perp,
\]

the required rank \(2n-2\) is obtained exactly in the alternatives (2.4)--(2.6).
Finally,

\[
\dim C(f)
=
(n-k)^2+k^2,
\tag{7.9}
\]

so under equality

\[
\dim(C(f)\cap C(g))
=
(n-k)^2+k^2-2n+3.
\tag{7.10}
\]

This completes the rank-one-side classification.

\begin{enumerate}
\item The exceptional dimension \(2\)
\end{enumerate}
Let \(n=2\) and let \(f\) be non-scalar. There is \(v\) such that \(v,fv\) are linearly independent; otherwise every vector would be an eigenvector, forcing \(f\) to be scalar.
An endomorphism commuting with \(f\) is determined by its value on \(v\), because

\[
X(fv)=f(Xv).
\]

Therefore

\[
C(f)=K[I,f],
\qquad
\dim C(f)=2.
\]

The same holds for \(g\). Both commutants contain \(KI\), and therefore

\[
\dim(C(f)+C(g))=3
\]

if and only if their intersection is exactly \(KI\), equivalently if and only if the two commutants are different.
Moreover,

\[
fg=gf
\iff g\in C(f)
\iff K[I,g]\subseteq C(f)
\iff C(g)=C(f).
\]

Hence

\[
\boxed{
\dim(C(f)+C(g))=3\iff fg\ne gf.
}
\]

This also explains why the descent argument in Lemma 4.1 cannot be imposed in dimension \(2\): a non-scalar \(2\times2\) matrix with irreducible quadratic minimal polynomial need not differ from a \(K\)-scalar by a \(K\)-rank-one map.

\begin{enumerate}
\item Attainment of the maximum
\end{enumerate}
Choose a basis \(e_1,\ldots,e_n\), with dual basis
\(\varepsilon_1,\ldots,\varepsilon_n\), and set

\[
f=e_1\otimes\varepsilon_1,
\qquad
g=(e_1+e_2)\otimes(\varepsilon_1+\varepsilon_2).
\tag{9.1}
\]

For \(f\), the eigenspaces have dimensions \(1\) and \(n-1\). Both the vector and covector defining \(g\) have nonzero components in both eigenspaces. Hence the split conditions (2.2) hold, and

\[
\dim(C(f)+C(g))=n^2-1.
\]

This construction is valid over every field. In characteristic \(2\),

\[
g^2=0,
\]

but \(g\) remains a nonzero rank-one endomorphism and the dimension calculation is unchanged.
Together with Theorem 3.1, this proves the exact maximum.

\begin{enumerate}
\item Dimension \(6\) and characteristic \(2\)
\end{enumerate}
The new bound immediately gives, over every field,

\[
\boxed{
n=6,\quad
\text{neither member scalar-plus-rank-one}
\Longrightarrow
\dim(C(f)+C(g))\le34.
}
\tag{10.1}
\]

Thus an extremal value \(35\) necessarily has a rank-one member. There is no exceptional characteristic-\(2\) family.
As an independent finite-field audit, I exhaustively checked the formerly possible \(\mathbf F_2\) obstruction. In the old dimension-count reduction, a counterexample without a rank-one member would have required

\[
\dim C(f)=\dim C(g)=20,
\qquad
\dim(C(f)\cap C(g))=5.
\]

Up to scalar shifts and similarity, the two relevant centralizer types are represented by

\[
P=\operatorname{diag}(1,1,0,0,0,0)
\]

and the rank-two square-zero map

\[
N(e_3)=e_1,\qquad N(e_4)=e_2,
\qquad N(e_i)=0\ \text{otherwise}.
\]

The exact enumeration covered:

\[
166\,656
\]

rank-two projections and

\[
136\,710
\]

rank-two square-zero maps. For every one of the four combinations

\[
(P,P),\quad(P,N),\quad(N,P),\quad(N,N),
\]

the minimum common-centralizer dimension was

\[
\boxed{6},
\]

never \(5\). The full histograms and exact orbit certificates are in the attached output. This computation is not needed for the proof above, but independently confirms the first unresolved case.


Exact C++ verifier


Exact enumeration output


SHA-256 checksums



\begin{enumerate}
\item Concise answer to Q655
\end{enumerate}
For every arbitrary field \(K\) and every \(n\ge2\),

\[
\boxed{\max\dim(C(f)+C(g))=n^2-1.}
\]

For \(n=2\), equality holds precisely when \(f\) and \(g\) do not commute.
For \(n\ge3\), equality holds precisely when, after interchanging \(f,g\), one member is a scalar plus a nonzero rank-one map and the other has quadratic minimal polynomial satisfying the split, irreducible, or square-zero conditions in Sections 2 and 7. There are no further radical, Jordan-module, finite-field, or characteristic-\(2\) exceptions.
Verification checklist


Canonical Q655 statement checked directly in the PDF: yes.


Exact maximum \(n^2-1\) over every field: proved.


Arbitrary-dimensional rank-one necessity: proved by Theorem 3.1.


Complete equality classification for all \(n\): proved.


Dimension \(6\): settled over every field.


Characteristic \(2\), including inseparable quadratics: covered.


Exact \(\mathbf F_2\) obstruction search: exhaustive, reproducible, no counterexample.


All commutant and intersection dimensions in the equality cases: given explicitly.
